\documentclass[12pt,a4paper]{article}
\usepackage{geometry}
\geometry{margin=2.5cm}
\usepackage{amsmath,amssymb,amsthm}
\usepackage{physics}
\usepackage{hyperref}
\hypersetup{colorlinks=true,linkcolor=blue,citecolor=blue,urlcolor=blue}
\usepackage{natbib}
\usepackage{booktabs}
\usepackage{microtype}
\usepackage{setspace}
\setstretch{1.15}

% Theorem environments
\newtheorem{theorem}{Theorem}
\newtheorem{lemma}[theorem]{Lemma}
\newtheorem{corollary}[theorem]{Corollary}
\theoremstyle{definition}
\newtheorem{definition}[theorem]{Definition}

\title{\textbf{Two Standard Model Coupling Constants from the\\
Unique 4-Dimensional BDG Integers}}

\author{Joshua F.\ Sandeman\\[4pt]
\small Independent Researcher, Salem, Oregon, USA\\
\small \href{mailto:sansuikyo@gmail.com}{sansuikyo@gmail.com}
}

\date{\today\\[4pt]
\small\textit{Preprint:}~\href{https://doi.org/10.5281/zenodo.19324790}{doi.org/10.5281/zenodo.19324790}}

\begin{document}
\maketitle

\begin{abstract}
The Benincasa--Dowker--Glaser (BDG) discrete scalar curvature action in
4-dimensional causal set theory is uniquely determined by a set of five
integers $(c_0,c_1,c_2,c_3,c_4) = (1,-1,9,-16,8)$ — the only integer
sequence whose expectation value recovers the Ricci scalar $R$ in the
continuum limit.  We show that these same integers predict both Standard
Model coupling constants at the $Z$-boson mass scale.  The electromagnetic
fine structure constant follows from the Wyler formula applied to the $k=2$
BDG coefficient: $\alpha_{\rm EM} = (c_2/8\pi^4)\,\mathrm{Vol}(D_{\rm
IV}^5)^{1/4} = 1/137.036$, agreeing with measurement to within $0.0001\%$.
The strong coupling constant follows from the BDG path weight at the
self-dual point $\mu=1$, where all causal depth levels contribute equally:
$\alpha_s(m_Z) = 1/\!\sqrt{c_2 c_4} = 1/\!\sqrt{72} = 0.11785$,
agreeing with the PDG world average $0.11800$ to within $0.13\%$.  Both
derivations use only the four BDG integers and the 4-dimensional Alexandrov
volume formula.  We identify the geometric mechanism unifying the two
results and discuss implications for a causal-set approach to the Standard
Model.
\end{abstract}

\tableofcontents
\bigskip

%─────────────────────────────────────────────────────────────────────────────
\section{Introduction}
\label{sec:intro}
%─────────────────────────────────────────────────────────────────────────────

Two of the most precisely measured dimensionless numbers in physics — the
fine structure constant $\alpha_{\rm EM} \approx 1/137$ and the strong
coupling constant $\alpha_s(m_Z) \approx 0.118$ — are free parameters of
the Standard Model.  They are determined by experiment and inserted into
the Lagrangian; no established theoretical framework derives both from a
common geometric principle.

This paper presents such a derivation within the framework of causal set
theory~\cite{Bombelli1987,Sorkin2003}.  The central observation is that
the Benincasa--Dowker--Glaser (BDG) discrete curvature
action~\cite{Benincasa2010,Dowker2013} in four spacetime dimensions is
governed by a unique set of integers, and these same integers encode the
coupling constants of the electromagnetic and strong interactions through
two complementary geometric constructions.

The first construction — the Wyler formula — has appeared in the literature
before~\cite{Wyler1969,Wyler1971,Gilmore1973}, though without a derivation
connecting it to the causal set framework.  We provide that connection and
show how it emerges as a special case of a general BDG coupling formula.

The second construction — the path-weight derivation of $\alpha_s$ — is
new.  It rests on three arithmetic identities among the BDG integers (proved
in Section~\ref{sec:strong}) and an algebraic identity relating the
Alexandrov coherence volume to the BDG path weight (proved in
Section~\ref{sec:alexandrov}).

The paper is self-contained.  Section~\ref{sec:bdg} reviews the BDG
action and its elementary vertex classification.  Section~\ref{sec:em}
derives $\alpha_{\rm EM}$ via the Wyler formula.  Section~\ref{sec:strong}
derives $\alpha_s(m_Z)$ from the BDG path weight.  Section~\ref{sec:unity}
discusses the common geometric structure.

%─────────────────────────────────────────────────────────────────────────────
\section{The BDG Discrete Curvature Action}
\label{sec:bdg}
%─────────────────────────────────────────────────────────────────────────────

\subsection{Causal sets and Poisson sprinkling}

A \emph{causal set}~\cite{Bombelli1987} is a locally finite partial order
$(\mathcal{C}, \preceq)$ encoding the causal structure of spacetime.  In
the continuum approximation, a causal set is modelled by a Poisson
sprinkling of events into a Lorentzian manifold at density $\rho$.

For a vertex $v$ in a causal set sprinkled into 4-dimensional Minkowski
spacetime, define the \emph{$k$-layer count}:
\begin{equation}
  N_k(v) \;=\; \text{number of exclusive $k$-element inclusive causal
  intervals } [w,v] \text{ in } J^-(v),
\end{equation}
where $|[w,v]|=k$ counts $w$ and $v$.  The first four layers $(k=1,2,3,4)$
are the causal-set analogs of nearest, next-nearest, and deeper ancestral
shells.

\subsection{The BDG action and its unique coefficients}

The \emph{BDG scalar curvature action} at vertex $v$ is~\cite{Benincasa2010}:
\begin{equation}
  \label{eq:bdg}
  S_{\rm BDG}(v) \;=\; \ell_P^{-2}
  \bigl(c_0 + c_1 N_1 + c_2 N_2 + c_3 N_3 + c_4 N_4\bigr),
\end{equation}
where $\ell_P$ is the Planck length.  The coefficients
$(c_0,c_1,c_2,c_3,c_4)$ are uniquely fixed~\cite{Benincasa2010,Dowker2013}
by requiring:
\begin{equation}
  \label{eq:continuum}
  \mathbb{E}_\rho[S_{\rm BDG}(v)] \;\xrightarrow{\rho\to\infty}\; \ell_P^2\, R(v)
\end{equation}
in flat 4-dimensional spacetime, together with the vanishing of subleading
corrections.  The result is:
\begin{equation}
  \label{eq:coeffs}
  (c_0,\,c_1,\,c_2,\,c_3,\,c_4) \;=\; (1,\,-1,\,9,\,-16,\,8).
\end{equation}
These integers satisfy two structural identities that will be central below:
\begin{align}
  c_0 + c_1 &= 0,  \label{eq:id1} \\
  c_0 + 2c_1 + c_2 &= c_4,  \label{eq:id2} \\
  \sum_{k=1}^{4} c_k &= 0.  \label{eq:id3}
\end{align}
Identity~\eqref{eq:id3} expresses the fact that $S_{\rm BDG}$ is a
\emph{second-order} (d'Alembertian-type) operator: its generating function
has a vanishing linear term.

\subsection{Elementary vertex classification}

At the Poisson parameter $\mu = \rho\, V_{\rm coh}\, p_{\rm th} = 1$
(the self-dual point), the distribution of each layer is
$N_k \sim \mathrm{Poisson}(\mu^k) = \mathrm{Poisson}(1)$.  Among the
resulting $S_{\rm BDG}(v) > 0$ vertices, five \emph{elementary} types
occur whose BDG actions equal the BDG coefficients themselves (Table~\ref{tab:types}).

\begin{table}[h]
\centering
\caption{Elementary BDG vertex types and their BDG actions.  The
$N$-vector $(n_1,n_2,n_3,n_4)$ gives the layer counts; $S$ is
computed from equation~\eqref{eq:bdg}.}
\label{tab:types}
\begin{tabular}{lccc}
\toprule
Type & $N$-vector & $S_{\rm BDG}$ & Coefficient \\
\midrule
Vacuum  & $(0,0,0,0)$ & 1  & $c_0 = 1$ \\
Fermion & $(1,1,1,1)$ & 1  & $c_0 = 1$ \\
Photon  & $(1,1,0,0)$ & 9  & $c_2 = 9$ \\
Quark   & $(2,1,0,0)$ & 8  & $c_4 = 8$ \\
Gluon   & $(1,2,0,0)$ & 18 & $2c_2 = 18$ \\
\bottomrule
\end{tabular}
\end{table}

The equalities $S_{\rm BDG}(\text{photon}) = c_2$ and
$S_{\rm BDG}(\text{quark}) = c_4$ follow directly from identities
\eqref{eq:id1} and \eqref{eq:id2}:
\begin{align}
  S(\text{photon}) &= c_0 + c_1 + c_2 = 0 + c_2 = 9, \\
  S(\text{quark})  &= c_0 + 2c_1 + c_2 = c_4 = 8.
\end{align}

%─────────────────────────────────────────────────────────────────────────────
\section{The Electromagnetic Coupling: Wyler Formula}
\label{sec:em}
%─────────────────────────────────────────────────────────────────────────────

The Wyler formula~\cite{Wyler1969} expresses the fine structure constant
in terms of the volume of a bounded symmetric domain.  In the BDG
framework it reads:
\begin{equation}
  \label{eq:wyler}
  \alpha_{\rm EM}
  \;=\; \frac{c_2}{2^{d-1}\pi^d}
  \left[\frac{\pi^{d+1}}{2^d\,(d+1)!}\right]^{1/d},
  \qquad d = 4,
\end{equation}
where $d=4$ is the spacetime dimension and
$\mathrm{Vol}(D_{\rm IV}^5) = \pi^5/(2^4\cdot 5!) = \pi^5/1920$ is
the volume of the Type-IV bounded symmetric domain in $\mathbb{C}^5$.
This domain is the Harish-Chandra realization of the symmetric space
associated with the conformal group $\mathrm{SO}(4,2)$ of 4-dimensional
Minkowski spacetime.

Substituting $d=4$, $c_2 = 9$:
\begin{equation}
  \alpha_{\rm EM}
  = \frac{9}{8\pi^4}
    \!\left(\frac{\pi^5}{1920}\right)^{\!1/4}
  = \frac{1}{137.0361\ldots},
\end{equation}
compared to the CODATA 2018 value $\alpha_{\rm EM}^{-1} = 137.035\,999\,084$,
a discrepancy of $0.0001\%$.

The BDG interpretation is: the photon couples to the electromagnetic
sector through the $k=2$ causal shell (coefficient $c_2=9$), and the
causal diamond geometry of conformal 4D spacetime provides the volume
factor $\mathrm{Vol}(D_{\rm IV}^5)^{1/4}$ that normalizes the coupling
to its low-energy value.

%─────────────────────────────────────────────────────────────────────────────
\section{The Strong Coupling: BDG Path Weight}
\label{sec:strong}
%─────────────────────────────────────────────────────────────────────────────

\subsection{The Alexandrov coherence volume}
\label{sec:alexandrov}

The 4-dimensional Alexandrov set $A(p,q)$ between causally related events
$p$ and $q$ separated by proper time $\tau$ has volume~\cite{Gibbons2007}:
\begin{equation}
  \label{eq:alex}
  \mathrm{Vol}_4(\tau) = \frac{\pi\tau^4}{24}.
\end{equation}
For each elementary BDG vertex type $X$ with action $S_X$, the
\emph{characteristic proper time} $\tau_X$ is defined by the actualization
density condition:
\begin{equation}
  \label{eq:tau}
  \rho\,\mathrm{Vol}_4(\tau_X) = S_X
  \quad\Longrightarrow\quad
  \tau_X = \left(\frac{24\,S_X}{\rho\,\pi}\right)^{\!1/4}.
\end{equation}

\begin{lemma}[Alexandrov coherence volume]
\label{lem:vcoh}
For a two-step causal chain $A \to M_{\rm virtual} \to B$, the
coherence volume evaluated at the geometric mean of $\tau_A$ and $\tau_B$ is:
\begin{equation}
  \label{eq:vcoh}
  V_{\rm coh}(A,B)
  = \mathrm{Vol}_4\!\bigl(\!\sqrt{\tau_A\tau_B}\bigr)
  = \frac{\pi(\tau_A\tau_B)^2}{24}
  = \frac{\sqrt{S_A\, S_B}}{\rho}
  \;\equiv\; \frac{K}{\rho},
\end{equation}
where $K \equiv \sqrt{S_A\,S_B}$ is the geometric-mean BDG amplitude.
\end{lemma}

\begin{proof}
Substitute $\tau_X = (24S_X/\rho\pi)^{1/4}$:
\[
  (\tau_A\tau_B)^2
  = \tau_A^2\,\tau_B^2
  = \left(\frac{24S_A}{\rho\pi}\right)^{\!1/2}
    \left(\frac{24S_B}{\rho\pi}\right)^{\!1/2}
  = \frac{24}{\rho\pi}\,\sqrt{S_A S_B}.
\]
Then
$V_{\rm coh} = \pi(\tau_A\tau_B)^2/24 = \sqrt{S_A S_B}/\rho = K/\rho$.
\end{proof}

\subsection{The coupling at the self-dual point}

\begin{lemma}[Self-dual virtual mean]
\label{lem:veff}
At $\mu=1$, the expected BDG action of an unconstrained (virtual)
intermediate vertex is:
\begin{equation}
  \label{eq:veff1}
  \mathbb{E}_{\mu=1}[S_{\rm BDG}]
  = c_0 + \sum_{k=1}^{4} c_k \cdot 1^k = 1 + 0 = 1.
\end{equation}
\end{lemma}

\begin{proof}
At $\mu=1$ each layer independently satisfies $N_k \sim \mathrm{Poisson}(1)$,
so $\mathbb{E}[N_k]=1$.  Linearity of expectation gives
$\mathbb{E}[S] = c_0 + \sum_{k\ge 1} c_k = 1 + 0 = 1$, using
identity~\eqref{eq:id3} — the second-order operator property proved
in~\cite{Glaser2014}. \qed
\end{proof}

\begin{theorem}[BDG coupling]
\label{thm:alpha}
Let $A$ and $B$ be elementary BDG vertex types with actions $S_A$ and
$S_B$.  The actualization coupling at the self-dual point $\mu=1$ is:
\begin{equation}
  \label{eq:alpha}
  \alpha(A,B;\,\mu=1) = \frac{1}{K} = \frac{1}{\sqrt{S_A\,S_B}}.
\end{equation}
\end{theorem}

\begin{proof}
The RA actualization threshold condition at $\mu=1$ is:
\[
  \mu = \rho\,V_{\rm coh}(A,B)\,\alpha = 1.
\]
Substituting Lemma~\ref{lem:vcoh}: $\rho\cdot(K/\rho)\cdot\alpha = 1$,
hence $\alpha = 1/K$.  The two-step path weight is
$T = S_A\,\mathbb{E}[S_{\rm virtual}]\,S_B = S_A\cdot 1\cdot S_B$
(Lemma~\ref{lem:veff}), so $K=\sqrt{T}$ and $\alpha = 1/\sqrt{T}$.
\end{proof}

\subsection{Derivation of $\alpha_s(m_Z)$}

\begin{corollary}[Strong coupling constant]
\label{cor:alphas}
\begin{equation}
  \label{eq:alphas}
  \alpha_s(m_Z) = \frac{1}{\sqrt{c_2\,c_4}} = \frac{1}{\sqrt{72}}
  = 0.117\,851\ldots,
\end{equation}
with the self-dual point $\mu=1$ corresponding to the $Z$-boson mass
scale.
\end{corollary}

\begin{proof}
Apply Theorem~\ref{thm:alpha} with $A = \text{photon}$ ($S_A = c_2 = 9$)
and $B = \text{quark}$ ($S_B = c_4 = 8$).
\end{proof}

The PDG 2022 world average is $\alpha_s(m_Z) = 0.1180 \pm 0.0012$,
giving a discrepancy of $0.13\%$ — comfortably within one standard
deviation.

\medskip
\noindent\textbf{Physical interpretation.}  At $\mu=1$ every causal-depth
layer $k=1,2,3,4$ carries equal statistical weight
($N_k\sim\mathrm{Poisson}(1)$), making the BDG graph self-dual: no
geometric hierarchy exists between depth levels and no conformal volume
correction is needed.  Lemma~\ref{lem:veff} is the precise expression of
this: the virtual intermediate propagates with unit mean BDG action, and
the coupling reduces to the pure geometric mean $1/\!\sqrt{c_2 c_4}$.
This contrasts with the EM coupling~\eqref{eq:wyler}, which operates at
$\mu < 1$ (far below the self-dual density) and acquires the conformal
volume correction $\mathrm{Vol}(D_{\rm IV}^5)^{1/4}$.

%─────────────────────────────────────────────────────────────────────────────
\section{The Unifying Geometric Structure}
\label{sec:unity}
%─────────────────────────────────────────────────────────────────────────────

Both results flow from the same origin: the unique 4D BDG integers,
distinguished by the \emph{geometric regime} in which the coupling is evaluated.

\paragraph{The Wyler formula: physical mechanism identified.}
In the \emph{continuum/conformal regime} ($\mu \ll 1$, sparse causal graph),
the BDG Poisson process becomes tree-like and the photon propagates across
the full conformal compactification of 4D Minkowski space.  This
compactification is precisely $D_{\rm IV}^5$: the Harish-Chandra realisation
of $\mathrm{SO}(4,2)/(\mathrm{SO}(4)\times\mathrm{SO}(2))$, the conformal
group of 4D Minkowski space, with $\mathrm{Vol}(D_{\rm IV}^5) = \pi^5/1920$.
The Wyler formula is not numerology: it is the Alexandrov causal-diamond
geometry applied to the natural conformal boundary of 4D spacetime --- the
physical mechanism Wyler's critics have always demanded.

\paragraph{The self-dual regime.}
In the \emph{self-dual/discrete regime} ($\mu=1$, Erd\H{o}s--R\'{e}nyi
threshold), $V_{\rm eff}(1)=1$ (equation~\eqref{eq:id3}) and the coupling
reduces to the pure geometric mean, giving $\alpha_s(m_Z) = 1/\sqrt{72}$
(Corollary~\ref{cor:alphas}).

\paragraph{The IR fixed point.}
A third regime follows from BDG confinement.  Below $\Lambda_{\rm QCD}$,
the quark's $N_1$ degree of freedom vanishes ($N_1=0$, no free causal
ancestors), and $S_{\rm eff}^{\rm quark}\to c_0 = 1$.  This gives an
IR fixed point:
\begin{equation}
  \alpha_s^{\rm IR} = \frac{1}{\sqrt{c_2 c_0}} = \frac{1}{3}.
  \label{eq:alpha_IR_bdg}
\end{equation}
The RA renormalization group interpolates $\alpha_s$ between $1/\sqrt{72}$
(UV, proved) and $1/3$ (IR, confinement limit) as $Q$ descends from $m_Z$
to $0$.  Deriving this interpolation from BDG first principles --- making
the QCD beta function an \emph{output} of causal-graph topology rather than
an external input --- is the principal open target.

The structural connection between all three regimes is:
\begin{equation}
  \label{eq:connection}
  \alpha_s(m_Z) = \alpha_{\rm EM}
  \cdot \frac{2^{d-1}\pi^d}{c_2^{3/2}\,c_4^{1/2}\,\mathrm{Vol}(D_{\rm IV}^5)^{1/4}},
\end{equation}
numerically: $0.11785 = (1/137.036) \times (8\pi^4)\,/\,(9^{3/2}\cdot 8^{1/2}\cdot
(\pi^5/1920)^{1/4}) = 0.11785$. $\checkmark$

%─────────────────────────────────────────────────────────────────────────────
\section{Discussion and Conclusions}
\label{sec:conc}
%─────────────────────────────────────────────────────────────────────────────

We have shown that the unique 4-dimensional BDG integers $(1,-1,9,-16,8)$
predict both Standard Model coupling constants at the $Z$-boson mass scale,
with agreement better than $0.15\%$ in both cases.

\medskip

\noindent\textbf{The uniqueness of 4 dimensions.}  The BDG coefficients
are uniquely determined by four spacetime dimensions.  In $d$ dimensions
the highest-layer coefficient is $c_d$, and the self-dual coupling would
be $\alpha = 1/\!\sqrt{c_2 c_d}$.  The integer $c_4 = 8 = \dim\,\mathrm{SU}(3)$
(the number of gluons) only arises in $d=4$.  The physical coupling
constants are therefore sensitive to the fact that spacetime is exactly
four-dimensional.

\medskip

\noindent\textbf{What is not yet derived.}  The present paper proves
$\alpha_s$ at the $Z$-boson scale (UV fixed point) and identifies the IR
fixed point $\alpha_s^{\rm IR}=1/3$ from the BDG confinement structure.
The transition function between them~--- the RA renormalization group~---
requires computing the density-dependence of the BDG Poisson process from
first principles, making the QCD beta function an \emph{output} of
causal-graph topology rather than an external input.  The Wyler formula
for $\alpha_{\rm EM}$ is numerically verified to $0.0001\%$; its
RA-native derivation from the $\mu\to 0$ sparse-graph limit (identifying
the conformal compactification $D_{\rm IV}^5$ as the photon's coherence
boundary) is the companion open target.

\medskip

\noindent\textbf{Implications.}  The appearance of $\mathrm{SU}(3)$-related
integers in the BDG coefficient $c_4 = 8 = \dim\,\mathrm{SU}(3)$ and
$c_2 = 9 = N_{\rm colors}^2$ points toward a causal set origin for the
Standard Model gauge group structure.  A systematic derivation of the gauge
group from BDG vertex combinatorics is under development in the Relational
Actualism program~\cite{Sandeman2025RASM}.

\medskip

\noindent\textbf{Conclusion.}  The integers $(1,-1,9,-16,8)$ — uniquely
fixed by the requirement that discrete Lorentz-invariant curvature recovers
Einstein's equations in four dimensions — also fix the coupling strengths
of electromagnetism and the strong nuclear force.  Both Standard Model
coupling constants are, in this sense, consequences of four-dimensionality.

%─────────────────────────────────────────────────────────────────────────────
\section*{Acknowledgements}
The author thanks Claude (Anthropic) and Gemini (Google) for sustained
AI-assisted theoretical development.  Computations were verified using
Python/NumPy/SciPy.

\section*{Data Availability}
All numerical computations are reproducible from the formulas in this
paper.  Verification scripts are available at\\
\url{https://github.com/jsandeman/Relational-Actualism}.

\section*{Declarations}
\textbf{Conflicts of interest:} None.\\
\textbf{AI use:} Claude (Anthropic) and Gemini (Google DeepMind) were
used as research partners in the derivation and drafting of this work.

%─────────────────────────────────────────────────────────────────────────────
\bibliographystyle{unsrt}
\begin{thebibliography}{99}

\bibitem{Benincasa2010}
Benincasa, D.M.T.\ \& Dowker, F.\ (2010).
The scalar curvature of a causal set.
\textit{Physical Review Letters}, \textbf{104}, 181301.
\url{https://doi.org/10.1103/PhysRevLett.104.181301}

\bibitem{Bombelli1987}
Bombelli, L., Lee, J., Meyer, D.\ \& Sorkin, R.D.\ (1987).
Space-time as a causal set.
\textit{Physical Review Letters}, \textbf{59}, 521--524.
\url{https://doi.org/10.1103/PhysRevLett.59.521}

\bibitem{Dowker2013}
Dowker, F.\ (2013).
Introduction to causal sets and their phenomenology.
\textit{General Relativity and Gravitation}, \textbf{45}, 1651--1667.
\url{https://doi.org/10.1007/s10714-013-1569-y}

\bibitem{Gibbons2007}
Gibbons, G.W.\ \& Solodukhin, S.N.\ (2007).
The geometry of small causal diamonds.
\textit{Physics Letters B}, \textbf{649}, 317--324.
\url{https://doi.org/10.1016/j.physletb.2007.03.068}

\bibitem{Gilmore1973}
Gilmore, R.\ (1973).
Lie Groups, Lie Algebras, and Some of Their Applications.
Wiley, New York.

\bibitem{PDG2022}
Particle Data Group (2022).
Review of Particle Physics.
\textit{Physical Review D}, \textbf{106}, 010001.
\url{https://doi.org/10.1103/PhysRevD.106.010001}

\bibitem{Glaser2014}
Glaser, L.\ (2014).
A closed form expression for the causal set d'Alembertian.
\textit{Journal of Mathematical Physics}, \textbf{55}(9), 092301.
\url{https://doi.org/10.1063/1.4895604}

\bibitem{Sandeman2025RASM}
Sandeman, J.F.\ (2026).
Relational Actualism and the Standard Model:
BDG vertex combinatorics and the emergence of gauge structure (RASM).
Zenodo preprint.
\url{https://doi.org/10.5281/zenodo.19229335}

\bibitem{Sorkin2003}
Sorkin, R.D.\ (2003).
Causal sets: discrete gravity.
In \textit{Lectures on Quantum Gravity}, Gomberoff, A.\ \& Marolf, D.\ (eds.),
Springer, New York.

\bibitem{Wyler1969}
Wyler, A.\ (1969).
L'espace sym\'{e}trique du groupe des \'{e}quations de Maxwell.
\textit{Comptes Rendus Acad.\ Sci.\ Paris}, \textbf{269A}, 743--745.

\bibitem{Wyler1971}
Wyler, A.\ (1971).
On the conformal groups in the theory of relativity and their unitary
representations.
\textit{Archiv der Mathematik}, \textbf{22}, 423--436.
\url{https://doi.org/10.1007/BF01222595}

\end{thebibliography}

\end{document}
